\section{Theorie}
\label{sec:Theorie}

\subsection{Das mikroskopische Modell}
Die stationäre Schrödingergleichung lautet 
\begin{equation*}
    \hat{H} \psi = E \psi.
\end{equation*}
Hierbei ist $\hat{H}$ der Hamilton-Operator, $\psi$ die Wellenfunktion und $E$ sind die Energieeigenwerte.
Umgestellt ergibt dies
\begin{equation*}
    E \psi = -\frac{\hbar^2}{2m_e} \Delta \psi + V(\vec{r}) \psi,
\end{equation*}
wobei für das Wasserstoffatom 
\begin{equation*}
    V(r) = -\frac{e^2}{4 \pi \epsilon_0 r}
\end{equation*}
gilt. 
Hierbei ist $e$ die Elementarladung, $\epsilon_0$ die elektrische Feldkonstante, $m_e$ die Masse des Elektrons, $\Delta$ der Laplace-Operator und $r$ der Abstand vom Elektron zum Atomkern. 
Um die Gleichung für die Energieeigenwerte zu lösen wird der Seperationsansatz verwendet. 
Der Ansatz für die Wellenfunktion lautet
\begin{equation*}
     \psi(r, \Theta, \Phi) = R(r) \cdot Y(\Theta, \Phi) 
\end{equation*}
Der Winkelanteil $Y(\Theta, \Phi)$ beschreibt die räumliche Orientierung der Wellenfunktion. 
Die Lösungen dieser Differentialgleichung sind die sogenannten Kugelflächenfunktionen $Y_{lm}(\Theta, \Phi)$. 
Diese lauten 
\begin{equation*}
    Y_l^m(\Theta, \Phi) = \sqrt{\frac{2l + 1}{4\pi} \frac{(l - m)!}{(l + m)!}} P_l^m(\cos \Theta) \, \mathrm{e}^{\mathrm{i}m\Phi} .
\end{equation*}
Hierbei ist l die Bahndrehimpulszahl und m die magnetische Quantenzahl. 
Die zugeordneten Legendre-Polynome
\begin{equation*}
    P_l^m(x) = \frac{(-1)^m}{2^l l!} (1 - x^2)^{m/2} \frac{\mathrm{d}^{m+l}}{\mathrm{d}x^{m+l}} (x^2 - 1)^l
\end{equation*}
stecken in diesen Funktionen.
Der Radialanteil $R_{nl}$ beschreibt, wie sich die Aufenthaltswahrscheinlichkeit des Elektrons mit dem Abstand $r$ vom Atom verhält mit 
\begin{equation*}
R_{nl} = N_{n,l} \cdot \left( \frac{2r}{na_0} \right)^l \cdot e^{-\frac{r}{na_0}} \cdot L_{n-l-1}^{2l+1} \left( \frac{2r}{na_0} \right).
\end{equation*}
Hierbei ist 
\begin{equation*}
N_{n,l} = a_0^{-\frac{3}{2}} \cdot \frac{2}{n^2} \cdot \sqrt{\frac{(n-l-1)!}{[(n+l)!]^3}}
\end{equation*}
der Normierungsfaktor.
Hierbei lauten die zugehörigen Laguerre-Polynome
\begin{equation*}
L_{\bar{n}}^j (x) = \frac{e^x x^{-j}}{\bar{n}!} \frac{\mathrm{d}^{\bar{n}}}{\mathrm{d}x^{\bar{n}}} (e^{-x} x^{\bar{n}+j}).
\end{equation*}

Die dazugehörigen Energieeigenwerte sind durch 
\begin{equation*}
    E_n = -\frac{e^4 m_e}{32 \pi^2 \epsilon_0^2 \hbar^2} \cdot \frac{1}{n^2}
\end{equation*}
beschrieben.
Das $n$ ist hierbei die Hauptquantenzahl und das $\hbar$ ist das reduzierte Plancksche Wirkungsquantum.

\subsection{Das makroskopische Modell}
Im Folgenden wird das makroskopische Modell betrachtet hierbei werden zur Untersuchung des Kugelresonators und des Zylinderresonators akustische Schallwellen verwendet. 

Die Helmholtzgleichung lautet 
\begin{equation*}
    (\Delta + k^2)p = 0, 
\end{equation*}
das $\Delta$ ist der Laplace-Operator, das $k$ die Wellenzahl und das $p$ ist die ortsabhängige Amplitude des Schalldrucks. 
Diese Gleichung kann auch durch den Seperationsansatz in Winkel- und Radialanteil seperiert werden.
Hierbei sind die Lösungen der Bewegungsgleichung für die Winkelanteile dieselben wie die Kugelflächenfunktionen $Y_{lm}(\Theta, \Phi)$, da die Helmholtzgleichung und die Schrödingergleichung in diesem Fall dieselbe Form haben. 
Für den radialen Teil ergibt sich
\begin{equation*}
    \frac{\partial^2 f(r)}{\partial r^2} + \frac{2}{r} \frac{\partial f(r)}{\partial r} + \left( k^2 - \frac{l(l+1)}{r^2} \right) f(r) = 0
\end{equation*}
Stehende Wellen können sich nur bilden, wenn ein Vielfaches ihrer Wellenlänge in den Hohlraumresonator passt. Hierbei stößt die Welle bei $r = R_k$ auf eine Wand. 
Das Argument für die Besselfunktion lautet $k_{ln} \cdot R_k = \Chi_{ln}$, woraus sich für die Wellenzahl 
\begin{equation*}
    k_{ln} = \frac{\Chi_{ln}}{R_k}
\end{equation*}
ergibt. 
Mithilfe der Formel $\omega_{ln} = k_{ln} \cdot c$ und der Wellenzahl $k_{ln}$ ergibt sich 
\begin{equation*}
    \omega_{ln} = \frac{\Chi_{ln}}{R_k} \cdot c.
\end{equation*}
Damit folgt 
\begin{equation}
    f_{ln} = \frac{\Chi_{ln}c}{2\pi R_k}
    \label{eq:frequnz_kugel}
\end{equation}
für die Frequenz.

Für den Zylinderresonator folgt, dass sich die Wellenzahl zusammensetzt aus dem Radialanteil und dem z-Anteil, mit dem Satz des Pythagoras ergibt sich
\begin{equation*}
    k^2 = k_r^2+k_z^2, 
\end{equation*}
wobei $k_z = \frac{\pi l}{L}$ und $k_r = \frac{\eta_{mn}}{R_z}$ für die k's eingesetzt werden, wobei $\eta_{mn}$ die $m$-te Nullstelle der Besselfunktion $J_n(kr)$ ist. Damit folgt 
\begin{equation*}
    k^2 = \left( \frac{\eta_{mn}}{R_Z} \right)^2 + \left( \frac{\pi l}{L} \right)^2.
\end{equation*}
Mithilfe von $\omega=k \cdot c$ und $\omega = 2 \pi f_{mnl}$ folgt
\begin{equation}
    f_{mnl} = \frac{c}{2\pi} \sqrt{\left( \frac{\eta_{mn}}{R_Z} \right)^2 + \left( \frac{\pi l}{L} \right)^2} .
    \label{eq:frequenz_zylinder}
\end{equation}

\clearpage